% Problemas del capítulo 9. Generado por tools/problemas/gen.py; no editar a mano.
\section{Problemas}\label{sec:problemas-09}

Graduados: \(\star\) directo, \(\star\star\) para pensar, \(\star\star\star\) desafío. Las respuestas están en el apéndice~\ref{ap:respuestas}; conviene intentarlos antes de mirarlas.

\begin{problema}[\({\star}\)]\label{pb:09-1}
Escribí \(v=(5,5)\) en la base \(b_1=(3,1)\), \(b_2=(1,2)\).
\end{problema}

\begin{problema}[\({\star}{\star}\)]\label{pb:09-2}
En una obra, el eje \(x\) forma \(30^\circ\) con el este (en sentido antihorario) y el origen está en \((E,N)=(1000;\ 500)\)~m. Un punto tiene coordenadas de obra \((20,10)\)~m. ¿Cuáles son sus coordenadas \((E,N)\)?
\end{problema}

\begin{problema}[\({\star}{\star}\)]\label{pb:09-3}
Diagonalizá \(A=\begin{pmatrix}4&1\\2&3\end{pmatrix}\) y usá \(A=CDC^{-1}\) para calcular \(A^5\).
\end{problema}

\begin{problema}[\({\star}{\star}\)]\label{pb:09-4}
La reflexión en la recta que forma \(30^\circ\) con el eje \(x\), en una base alineada con esa recta y su perpendicular, es \(\mathrm{diag}(1,-1)\). Pasala a la base canónica.
\end{problema}

\begin{problema}[\({\star}{\star}{\star}\)]\label{pb:09-5}
La forma \(\omega=(2,1)\) está escrita en la base canónica. Escribila en la base \(b_1=(3,1)\), \(b_2=(1,2)\), calculando \(\omega(b_1)\) y \(\omega(b_2)\). Comprobá que \(\omega(v)\) da lo mismo en las dos bases, con \(v=(5,5)\) del problema 1.
\end{problema}
